\documentclass[10pt,a4paper]{article}

\usepackage[utf8]{inputenc}
\usepackage[T1]{fontenc}
\usepackage[spanish,es-tabla]{babel}

\usepackage{amsmath,amssymb,amsthm}
\usepackage{booktabs}
\usepackage{tikz}
\usetikzlibrary{automata,positioning,arrows}
\pagecolor{black!85!white}
\color{white}

\theoremstyle{definition}
\newtheorem{definicion}{Definición}[section]

\title{Hito 0 - Automatas finitos: enteros positivos}
\author{Jonathan Javier Pietracci}
\begin{document}

\maketitle

Creando mi compilador en C;

\section{Lenguaje y alfabeto}

Sea el alfabeto de digitos decimales
\[
\Sigma = \{0,1,2,3,4,5,6,7,8,9\}.
\]
El lenguaje de los enteros positivos en notación decimal es el conjunto de cadenas sobre $\Sigma$ que cumplen:

\begin{enumerate}
    \item el primer simbolo es un digito del $1$ al $9$
    \item los simbolos siguientes pueden ser cualquiera de los $10$ digitos.
\end{enumerate}

\section{Gramática formal}

\begin{align*}
    \text{Numero} &\to \text{PrimerDigito} \ \text{Resto} \\
    \text{Resto} &\to \varepsilon \mid \text{SiguienteDigito} \ \text{Resto} \\
    \text{PrimerDigito} &\to 1 \mid 2 \mid 3 \mid 4 \mid 5 \mid 6 \mid 7 \mid 8 \mid 9 \\
    \text{SiguienteDigito} &\to 0 \mid 1 \mid 2 \mid 3 \mid 4 \mid 5 \mid 6 \mid 7 \mid 8 \mid 9 \\
\end{align*}
donde $\varepsilon$ denota la cadena vacia.

\section{Autómata finito determinista (AFD)}

\begin{definicion}
    Un automata finito determinista es una quíntupla
    $M = (S, I, f, s_0, F)$ donde $S$ es un conjunto finito de estados,
    $I$ es el alfabeto de entrada, $f : S \times I \to S$ es la función de
    transición, $s_0 \in S$ es el estado inicial y $F \subseteq S$ es el
    conjunto de estados de aceptación.
\end{definicion}

Para nuestro lenguaje:
\[
    S = \{q_0, q_1, q_2\} \quad
    I = \{0,1,2,3,4,5,6,7,8,9\}, \quad
    s_0 = q_0, \quad
    F = \{q_1\}
\]
La función de transición $f$ es:
\begin{align*}
    f(q_0, d) &= q_1 \quad \text{si } d \in \{1,\dots,9\}, \\
    f(q_0, 0) &= q_2, \\
    f(q_1, d) &= q_1 \quad \text{para todo } d \in I, \\
    f(q_2, d) &= q_2 \quad \text{para todo } d \in I \quad \text{(estado de rechazo hundido)}. \\
\end{align*}

\begin{table}[h]
\centering
    \begin{tabular}{c|cccccccccc}
        \toprule
        Estado & 0 & 1 & 2 & 3 & 4 & 5 & 6 & 7 & 8 & 9 \\
        \midrule
        $q_0$ (inicial) & $q_2$ & $q_1$ & $q_1$ & $q_1$ & $q_1$ & $q_1$ & $q_1$ & $q_1$ & $q_1$ & $q_1$ \\
        $q_1$ (aceptación) & $q_1$ & $q_1$ & $q_1$ & $q_1$ & $q_1$ & $q_1$ & $q_1$ & $q_1$ & $q_1$ & $q_1$ \\
        $q_2$ (rechaza) & $q_2$ & $q_2$ & $q_2$ & $q_2$ & $q_2$ & $q_2$ & $q_2$ & $q_2$ & $q_2$ & $q_2$ \\
        \bottomrule
    \end{tabular}
    \caption{Función de transición $f$. Doble círculo = estado de aceptación.}

\end{table}

\subsection{Diagrama de estados}

\begin{center}
    \shorthandoff{<>}
    \begin{tikzpicture}[shorten >=1pt,node distance=2.5cm,on grid,auto,every state/.style={draw,circle,minimum size=1cm}]
        \node[state,initial]   (q0)               {$q_0$};
        \node[state,accepting] (q1) [right=of q0] {$q_1$};
        \node[state]           (q2) [below=of q0] {$q_2$};

        \path[->]
          (q0) edge [bend left]  node {1..9} (q1)
          (q0) edge              node [left] {0} (q2)
          (q1) edge [loop right] node {0..9} (q1)
          (q2) edge [loop left]  node {0..9} (q2);
    \end{tikzpicture}
    \shorthandon{<>}
\end{center}

\subsection{Verificación por simulación}

Recorremos la cadena símbolo a símbolo usando la transición extendida.

\begin{table}[h]

    \centering
    
    \begin{tabular}{c|c|c}
        \toprule
        Cadena & Recorrido de estados & Veredicto \\
        \midrule
        \texttt{7} & $q_0 \rightarrow {7} q_1$ & acepta \\
        \texttt{42} & $q_0 \rightarrow {4} q_1 \rightarrow {2} q_1$ & acepta \\
        \texttt{40} & $q_0 \rightarrow {4} q_1 \rightarrow {0} q_1$ & acepta \\
        \texttt{102} & $q_0 \rightarrow {1} q_1 \rightarrow {0} q_1 \rightarrow{2} q_1$ & acepta\\
        \texttt{0} & $q_0 \rightarrow{0} q_2$ & rechaza \\
        \texttt{007} & $q_0 \rightarrow {0} q_2 \rightarrow {0} q_2 \rightarrow {7} q_2$ & rechaza \\
        \bottomrule
    \end{tabular}

    \caption{El AFD reconoce exactamente los enteros positivos sin ceros a la izquierda.}

\end{table}

\section{Gramáticas formales y jerarquía de Chomsky}

Las gramáticas se clasifican en cuatro tipos según la forma de sus reglas:

\begin{enumerate}
    \item tipo 0: sin restricciones
    \item tipo 1: sensibles al contexto ($A$ entre símbolos vecinos).
    \item tipo 2: libres de contexto (un solo símbolo a la izquierda).
    \item tipo 3: regulares (reconocibles por un autómata finito).
\end{enumerate}

\section{Teoría de lenguajes formales}

Un lenguaje formal $L$ sobre un alfabeto $\Sigma$ es un conjunto de cadenas de $\Sigma$.
La cerradura de Kleene $\Sigma^*$ es el conjunto de todas las cadenas sobre $\Sigma$,
incluyendo la cadena vacía.

Nuestro lenguaje de enteros positivos es un subconjunto propio de
$\Sigma^*$:
\[
    L = \{ w \in \Sigma^* \mid \text{el primer símbolo de } w \text{ es } 1,\dots,9 \}.
\]
No coincide con $\Sigma^*$ porque $\Sigma^*$ contiene cadenas como \texttt{0}, \texttt{00}, \texttt{007}, que nuestro AFD rechaza.
Como $L$ es reconocido por un autómata finito, es un lenguaje regular.

\section{Inducción estructural}

Demostramos por inducción sobre la construcción de las cadenas
que todo elemento del lenguaje generado por nuestra gramática empieza 
con un dígito del $1$ al $9$

\begin{proof}
    \textbf{Caso base.} La cadena más corta se obtiene con
    $\text{Resto} \to \varepsilon$, lo que da un $\text{PrimerDigito}$,
    por ejemplo $1$. Esta cadena empieza con $1$, luego cumple la propiedad.

    \textbf{Paso inductivo.} Supongamos que una cadena $w$ cumple la propiedad (empieza con $1$ al $9$).
    Al aplicar $\text{Resto} \to \text{SiguienteDigito}\ \text{Resto}$ se agrega
    un dígito $d \in \{0,\dots,9\}$ al final de $w$. Como $d$ se agrega al final,
    el primer símbolo de $w$ no cambia; por tanto $wd$ también empieza con $1$ al $9$.
    
    Por inducción, toda cadena generada por la gramática cumple la propiedad.
\end{proof}

\section{Árboles y el árbol de sintaxis (AST)}

Un árbol es un grafo conexo y acíclico con un nodo raíz; cada nodo tiene un único
padre, salvo la raíz. Los nodos sin hijos son las hojas.

En el compilador, una expresión se representa como un árbol de sintaxis abstracta (AST).
La expresión $1 + 2 \times 3$ se esctructura así:

\begin{center}
    \shorthandoff{<>}
    \begin{tikzpicture}[level distance=1.2cm,
        every node/.style={draw,circle,minimum size=0.8cm}]
        \node {+}
            child{ node {1} }
            child{ node {$\times$}
                child{ node {2} }
                child{ node {3} } };
    \end{tikzpicture}
    \shorthandon{<>}
\end{center}

Las hojas son $1, 2, 3$. Al evaluar de las hojas hacia la raíz se
calcula primero $2 \times 3 = 6$ y luego $1 + 6 = 7$, respetando la precedencia.

\section{Tablas hash y la tabla de símbolos}

\subsection{Función de hash}

Una función de hash es una aplicación
\[
    h : U \longrightarrow \{0,1,\dots,m-1\}
\]
que asigna a cada clave del universo $U$ (por ejemplo, todas las cadenas posibles) un índice
en un arreglo de $m$ casillas. La función $h$ se \emph{define}; no es una construcción condicional:
a cada clave le corresponde un único índice por la fórmula elegida.

\subsection{Principio del palomar}

Si $|U|$ (la cantidad de claves posibles) es infinito y $m$ es finito,
entonces forzosamente existen al menos dos claves distintas con el mismo
índice. Esto no es un error de implementación: es consecuencia del principio del palomar
(si se ponen más palomares que casillas alguna casilla recibe al menos dos). Por tanto,
las \emph{colisiones} son inevitables y deben \emph{resolverse}, no evitarse.

\subsection{El módulo como herramienta matemática}

Para claves enteras se usa
\[
    h(k) = k \bmod m,
\]
donde $k \bmod m$ es el resto único de dividir $k$ entre $m$, y por
definición cumple $0 \le (k \bmod m) < m$. Así el índice siempre
cae en el rango del arreglo. En la práctica se elige $m$ primo para
que la distribución no descarte bits (si $m$ fuera potencia de 2, $k \bmod m$
usaría solo los últimos bits y agruparía claves con patrones similares).

\subsection{Claves que no son enteras}

Si la clave es un \emph{string} (como los identificadores del compilador),
el módulo no aplica directamente porque opera sobre enteros. Hay que convertir primero
el string en un entero que dependa de \emph{todos} sus caracteres y de su orden:

\[
    e = 0;\quad \text{para cada caracter } c:\ e = e \cdot 256 + \text{código}(c)
\]

Por ejemplo, \texttt{'ab'} (códigos 97 y 98) da $e = 97\cdot 256 + 98 = 24930$. Entonces
$h(\text{\ttfamily 'ab'} ) = 24930 \bmod m$. La misma cadena produce
siempre el mismo entero, y por tanto la misma casilla.

\subsection{Resolución de colisiones}

\begin{itemize}
    \item \textbf{Encadenamiento:} cada casilla contiene una lista enlazada. Si 
        una segunda clave cae en una casilla ocupada, se agrega
        a la lista. Pueden guardarse más de $m$ claves (el límite es la memoria).
    \item \textbf{Direccionamiento abierto:} cada casilla guarda una sola
        clave. Si la casilla $i$ está ocupada, se prueba $i+1$,
        $i+2$, \dots\ hasta hallar una libre. A lo sumo caben
        $m$ claves; la tabla se satura cuando las casillas están llenas.
\end{itemize}

\begin{center}
    \shorthandoff{<>}
    \begin{tikzpicture}[
        casilla/.style={draw,rectangle,minimum width=2.2cm,minimum height=0.9cm},
        nodo/.style={draw,rectangle,minimum width=2.2cm,minimum height=0.7cm}]
        \matrix [row sep=0.25cm] {
            \node[casilla] (c0) {0}; & \node[nodo] {$k_1$};\\
            \node[casilla] (c1) {1}; & \node[nodo] {$k_2$}; \node[nodo,right=0.1cm of c1] (n1) {$k_4$}; \\
            \node[casilla] (c2) {2}; & \\
            \node[casilla] (c3) {3}; & \node[nodo] {$k_3$}; \node[nodo,right=0.1cm of c3] (n3) {$k_4$}; \\ 
        };
        \draw[->] (c1) -- (n1);
        \draw[->] (c3) -- (n3);

    \end{tikzpicture}
    \shorthandon{<>}

\end{center}

El dibujo muestra encadenamiento: las casillas 1 y 3 tienen listas enlazadas ($k_5$ cuelga de $k_2$; $k_4$ cuelga de $k_3$).

\subsection{Conexión con el compilador}

La tabla hash será la \emph{tabla de símbolos}: mapea identificadores reconocidos
por el lexer (Hito 1) a su información. El lexer produce el string, se aplica la conversión a entero
y el módulo, y se resuelve la colisión con la estrategia elegida.

\section{Binario y Sistema Posicional}

Antes de hablar de "binario", necesitamos entender \texttt{qué es un sistema posicional}. Es una idea poderosa:
representar cantidades usando \texttt{posiciones} y \texttt{pesos}.

la definición formal dice:

\begin{definicion}

    Un entero en base b se representa como
    \[
       N = d_{k-1} \cdot b^{k-1} + d_{k-2} \cdot b^{k-2}
            + \cdots + d_1 \cdot b^1 + d_0 \cdot b^0
        \]
       donde $0 <= d_i < b$ para cada $i$.
    
\end{definicion}

\textbf{Ejemplo en base 10:} el número $347_{10}$ se descompone como:

\begin{center}
    \begin{tabular}{c|ccc}
    \toprule
    Posición & 2 & 1 & 0 \\
    \midrule
    Dígito $d_i$ & 3 & 4 & 7 \\
    Peso $b^i$ & $10^2 = 100$ & $10^1 = 10$ & $10^0 = 1$\\
    \bottomrule
    \end{tabular}
\end{center}

\[
347_{10} = 3 \cdot 100 + 4 \cdot 10 + 7 \cdot 1 = 347
\]

la base b define: (1) cuántos dígitos posibles hay (del 0 al b-1) 
y (2) cuánto 'vale' cada posición (potencias de $b$).

\subsection{Binario: base 2}

En base 2 solo existen dos dígitos: \{0, 1\}. Los pesos son potencias de 2:

\begin{center}
\begin{tabular}{c|cccccc}
    
    \toprule
    Posición & 5 & 4 & 3 & 2 & 1 & 0 \\
    \midrule
    Peso $2^i$ & 32 & 16 & 8 & 4 & 2 & 1 \\
    \bottomrule

\end{tabular}
\end{center}

Ejemplo: representar $21_{10}$ en binario.
Se descompone $21 = 16 + 4 + 1$:

\begin{center}
\begin{tabular}{c|ccccc}
    \toprule
    Posición & 4 & 3 & 2 & 1 & 0 \\
    \midrule
    Peso $2^i$ & 16 & 8 & 4 & 2 & 1\\
    Dígito & 1 & 0 & 1 & 0 & 1 \\
    \bottomrule
\end{tabular}
\end{center}

 \[
     21_{10} = 10101_{2}
 \]

\subsection{Conversión decimal a binario}

\begin{definicion}
    Para convertir un entero positivo n a binario, se aplican divisiones
    sucesivas por 2. Los restos, leídos de último a primero, forman los dígitos
    binarios (del LSB al MSB).
\end{definicion}

Ejemplo: $47_{10}$ a binario.

\begin{center}
\begin{tabular}{c|c|c}
    \toprule
    División & Cociente & Resto \\
    \midrule
    $47 \div 2$ & 23 & 1 \\
    $23 \div 2$ & 11 & 1 \\
    $11 \div 2$ & 5 & 1 \\
    $5 \div 2$ & 2 & 1 \\
    $2 \div 2$ & 1 & 0 \\
    $2 \div 1$ & 0 & 1 \\
    \bottomrule
\end{tabular}
\end{center}

Los restos de abajo hacia arriba: $47_{10} = 101111_{2}$.

\subsection{Conversión binario a decimal}

\begin{definicion}
    Para convertir $d_{k-1} \cdots d_1 d_0$ a decimal, se calcula
    $\sum_{i=0}^{k-1} d_i \cdot 2^i$.
\end{definicion}

Ejemplo: $110011_{2}$ a decimal.

\begin{center}
    \begin{tabular}{c|cccccc}
        \toprule
        Posición & 5 & 4 & 3 & 2 & 1 & 0 \\
        \midrule
        Dígito & 1 & 1 & 0 & 0 & 1 & 1 \\
        Peso $2^i$ & 32 & 16 & 8 & 4 & 2 & 1 \\
        \bottomrule
    \end{tabular}
\end{center}

\[
    32 + 16 + 0 + 0 + 2 + 1 = 51_{10}
\]
\subsection{Hexadecimal: base 16}

\begin{definicion}
    El sistema hexadecimal usa 16 dígitos: $\{0,1,2,3,4,5,6,7,8,9,A,B,C,D,E,F\}$.
    Cada dígito hex equivale a exactamente 4 bits.
\end{definicion}

\begin{center}
    \begin{tabular}{c|c|c}
        \toprule
        Hex & Binario & Decimal \\
        \midrule
        0 & 0000 & 0 \\
        1 & 0001 & 1 \\
        2 & 0010 & 2 \\
        3 & 0011 & 3 \\
        4 & 0100 & 4 \\
        5 & 0101 & 5 \\
        6 & 0110 & 6 \\
        7 & 0111 & 7 \\
        8 & 1000 & 8 \\
        9 & 1001 & 9 \\
        A & 1010 & 10 \\
        B & 1011 & 11 \\
        C & 1100 & 12 \\
        D & 1101 & 13 \\
        E & 1110 & 14 \\
        F & 1111 & 15 \\
        \bottomrule
    \end{tabular}
\end{center}

Binario a hex: separar en nibbles (grupos de 4 bits)
y reemplazar cada grupo.

$11010101_{2} = \underbrace{1101}_{D}\;\underbrace{0101}_{5} = D5_{16}$

Hex a binario: cada dígito hex -> 4 bits.

$3F_{16} = \underbrace{3}_{0011}\;\underbrace{F}_{1111} = 00111111_{2}$

\subsection{Notación de bases}

Para desambiguar la base de un número se usa subíndice:
$101_2$ (cinco en binario), $101_{10}$ (ciento uno en decimal).

En C++ se usa prefijo: \texttt{0b101} (binario),
\texttt{0xD5} (Hexadecimal), \texttt{101} (decimal).

\section{Aritmética binaria}

\subsection{Suma binaria}

La suma binaria sigue las mismas reglas que la suma decimal, con la
salvedad de que 1 + 1 = 10 (cero con acarreo de 1).

\begin{center}
    \begin{tabular}{c|c|c|c}
        \toprule
        $A$ & $B$ & Resultado & Carry \\
        \midrule
        0 & 0 & 0 & 0 \\
        0 & 1 & 1 & 0 \\
        1 & 0 & 1 & 0 \\
        1 & 1 & 0 & 1 \\
        \bottomrule
    \end{tabular}
\end{center}

Se suma de derecha a izquierda (LSB $\to$ MSB), propagando el carry.
Si al final queda un carry, se escribe como bit adicional.

Ejemplo: $1101_{2}$ + $1011_{2}$

\begin{center}
\begin{tabular}{ccccc}
    \toprule
        & 1 & 1 & 1 &  \\
      + & 1 & 0 & 1 & 1 \\
        \hline
    1 & 1 & 0 & 0 & 0 \\
\end{tabular}
\end{center}

$13 + 11 = 24 = 11000_{2}$ \checkmark

\subsection{Complemento a 2}

\begin{definicion}
    El complemento a 2 de un número binario de n bits se obtiene:
    \begin{enumerate}
        \item Invertir todos los bits (complemento a 1).
        \item Sumar 1 al resultado.
    \end{enumerate}
    \[
        -x = \sim x + 1
    \]
\end{definicion}

Con $n$ bits, el rango de enteros con signo es
$[-2^{n-1},\; 2^{n-1} - 1]$.

Ejemplo -5 con 4 bits.

\begin{enumerate}
    \item 5 = 0101
    \item Invertir: 1010
    \item Sumar 1: 1011
\end{enumerate}

$-5_{10} = 1011_2$ \checkmark

\subsection{Tabla completa de 4 bits}

\begin{center}
\begin{tabular}{c|c||c|c}
    \toprule
    Decimal & Binario & Decimal & Binario \\
    \midrule
    0 & 0000 & -1 & 1111 \\
    1 & 0001 & -2 & 1110 \\
    2 & 0010 & -3 & 1101 \\
    3 & 0011 & -4 & 1100 \\
    4 & 0100 & -5 & 1011 \\
    5 & 0101 & -6 & 1010 \\
    6 & 0110 & -7 & 1001 \\
    7 & 0111 & -8 & 1000 \\
    \bottomrule
\end{tabular}
\end{center}

\subsection{Resta por suma}

La resta se implementa como suma con el negativo:
\[
    A - B = A + (-B)
\]
El carry que excede el ancho de bits se descarta. 

Ejemplo: $6 - 4$ con 4 bits.
$6 = 0110$, $-4 = 1100$ (invertir $0100$ $\to$ $1011$, sumar 1 $\to$ $1100$).
\[
    0110 + 1100 = 10010 \to 0010 = 2_{10} \quad \checkmark
\]
\subsection{Overflow}

\begin{definicion}
    Hay overflow signed cuando el carry de entrada al MSB es distinto
    del carry de salida del MSB. El resultado no cabe en el rango
    $[-2^{n-1},\; 2^{n-1} - 1]$.
\end{definicion}

Ejemplo: 7 + 1 = 8 con 4 bits (rango $[-8, 7]$).

$0111 + 0001 = 1000$. El MSB indica negativo (-8), pero la 
suma real es +8. Hay overflow.

\subsection{Un solo cero}

En complemento a 2, $-0 = 0000_2$ (mismo que $+0$).
No existe el doble cero, a diferencia del complemento a 1.

\end{document}
